\section{Discussion}
\label{sec:discussion}
\subsection{What the Results Establish}
The useful output is a verified decision about the exact quantized cost
change. An approximate predictor can agree with exact labels yet leave that
decision unresolved; conversely, an aggregate interval can resolve the
update without certifying every individual assignment. The evaluation
therefore separates prediction agreement, label certificates and update
coverage instead of treating them as one accuracy score.

Exact old membership and authenticated group moments remove uncertainty
from the assigned-center contribution. The rounded-basis envelope confines
the remaining uncertainty to switching, and selected original-row proofs
can remove part of that interval. The 19-update diagnostic and naturally
worsening archived Fashion requests illustrate these different roles.
Their practical value is measured through complete proof work, including
inconclusive certificates and exact fallback.

\subsection{Representation Choice and Computational Trade-offs}
The centroid-informed prefix is designed to preserve directions relevant
to the deployed centers. Training-residual directions extend that span,
while the certificate-only extension leaves prediction rank unchanged.
Future candidates can move outside the represented directions, increasing
residual uncertainty. The low-rank native contrast is consistent with this
mechanism, but does not establish a general runtime ranking between bases.

At rank 79, PCA is often close to the centroid-informed configuration. Both
instantiate the same verifier-enforced mechanism. The evidence supports the
certificate framework more clearly than a claim that the fitted centroid
basis is always the better predictor. A fair comparison of scheduling must
also optimize PCA; a routing improvement is not automatically a geometric
advantage over that basis.

The later gate provides that equal-policy comparison: PCA and centroid
resolve the same image requests, and PCA has a slightly lower warm total
in the single measured execution. Raw projections, coresets and RaBitQ
extend the numerical controls, but do not have interchangeable outputs.
In particular, better label agreement does not imply better cost-change
sign estimation, and good coreset estimates do not become authenticated
full-catalog decisions. The cold ZKML-adapter wins likewise prevent a
general superiority claim for the compact service.

Preparation is not dominated by the numerical dimension reduction in the
reported small catalogs. Authentication and key work account for much more
time, and recurring key loading matters even when setup is reusable. This
explains why cold, prepared and resident-service settings need different
measurements. It also explains why larger request counts do not guarantee
success when most requests still require expensive exact evaluation.

\subsection{Implications for Prototype-Update Audits}
A useful deployment compares multiple candidates with one committed catalog
and deployed reference. In that setting, the Fashion results show a measured
opportunity to avoid repeated original-space proof work. A known future
bank, a short reference lifetime or a workload with weak bounds can favor
other contracts. The text results are an operating limit of the compact
path and should remain visible when selecting the service policy.

The certificates establish an ordering of quantized distortion, not an
improvement in class accuracy or proof of encoder training. Public assigned
aggregates and corrections are also meaningful disclosures. The
moment-recovery construction in Section~\ref{sec:methodology-guarantees}
shows why zero knowledge relative to these outputs does not provide
adaptive record non-reconstruction.

\subsection{Remaining Evaluation}
The next bounded study prioritizes completed comparators, independently
fitted configurations and repeated paired timings. The later gate supplies
an equally optimized PCA comparison, but only one proof execution per cell.
Remaining priorities include complete-response preparation/request costs,
real-proof low-rank ablations and common-contract scale measurements.
Per-response timestamps must reflect a complete ready
answer, including any batch preflight and replay. Longer prefixes will be
reported only when their request proofs have actually been generated.

Reference-reset experiments will charge new authentication and any required
setup rather than assuming permanent reuse. Additional backends and an
exact sign-only control remain useful directions. These measurements are
not represented by estimated server times, unexecuted grid cells or the
single-execution ratios already available.

